%=============================================================================!
% 3.2  KINETIC ENERGY                                                         !
%=============================================================================!
\section{Kinetic energy}
\label{sec:en-kinetic}

The work done on a body must go somewhere. This section shows, from the
second law, that the work done by the net force changes one particular
property of the motion, a combination of mass and speed called the
kinetic energy.

\subsection*{Along a line}

Let a constant net force $F$ act along a line on a body of mass $m$. By
the second law its acceleration $a = F/m$ is constant, and
\cref{eq:mo-v2} relates its velocity $v$ at position $x$ to its velocity
$v_0$ at position $x_0$, $v^2 = v_0^2 + 2a(x - x_0)$. Multiplying both
sides by $m/2$,
\begin{align*}
   \tfrac12 m v^2 &= \tfrac12 m v_0^2 + m a\,(x - x_0)\\
   &= \tfrac12 m v_0^2 + F\,(x - x_0)
      && \text{by the second law, $ma = F$.}
\end{align*}
The last term is the work done by the net force over the displacement $x
- x_0$, so the work has been added to the quantity $\frac12 m v^2$. This
quantity is the kinetic energy of the body,
\begin{equation}
   K = \tfrac12 m v^2 ,
   \label{eq:en-kinetic}
\end{equation}
and the result reads
\begin{equation*}
   W = K - K_0 ,
\end{equation*}
where $K_0 = \frac12 m v_0^2$: the work done by the net force equals the
change in the kinetic energy. Its unit is the joule, as it must be, since
a kilogram times a speed squared is
\si{\kilogram\metre\squared\per\second\squared}.

The trolley of \cref{sec:ne-second}, of mass \SI{20}{\kilogram}, pushed
from rest by a net force of \SI{10}{\newton} through \SI{1}{\metre}, gains
$W = 10\times1 = \SI{10}{\joule}$ of kinetic energy. Setting $\frac12 m
v^2 = K$ and solving for the speed,
\begin{equation*}
   v = \sqrt{\frac{2K}{m}} = \sqrt{\frac{2\times10}{20}}
   = \SI{1}{\metre\per\second} .
\end{equation*}
The second law gives the same result the long way: the acceleration is
$10/20 = \SI{0.5}{\metre\per\second\squared}$, the trolley covers
\SI{1}{\metre} in \SI{2}{\second}, since $\frac12\times0.5\times2^2 = 1$,
and by then it moves at $0.5\times2 = \SI{1}{\metre\per\second}$. The
energy argument never needed the time.

\subsection*{In space}

The same steps work for a constant force in three dimensions. The
acceleration $\vb{a} = \vb{F}/m$ is constant, so each component obeys
\cref{eq:mo-v2} on its own, $v_x^2 = v_{0x}^2 + 2a_x\,\Delta x$, and
likewise for $y$ and $z$ (a component with no acceleration keeps a
constant velocity, so the equation holds for it too, although
\cref{eq:mo-v2} was derived for $a \neq 0$), where $\Delta\vb{r} = (\Delta x, \Delta y,
\Delta z)$ is the displacement. Adding the three equations, and writing
$v^2 = \vb{v}\cdot\vb{v} = v_x^2 + v_y^2 + v_z^2$ for the speed squared,
\begin{equation*}
   v^2 = v_0^2 + 2\,(a_x\,\Delta x + a_y\,\Delta y + a_z\,\Delta z)
   = v_0^2 + 2\,\vb{a}\cdot\Delta\vb{r} ,
\end{equation*}
and multiplying by $m/2$ and using $m\vb{a} = \vb{F}$,
\begin{equation*}
   \tfrac12 m v^2 - \tfrac12 m v_0^2 = \vb{F}\cdot\Delta\vb{r} = W .
\end{equation*}
The kinetic energy $\frac12 m v^2$ depends on the speed alone, not on the
direction of motion, and the work is that of \cref{eq:en-work-constant},
with $\Delta\vb{r}$ the displacement from start to finish even when the
path between them curves. The component of the force perpendicular to
$\Delta\vb{r}$ has indeed changed nothing that $K$ measures.

A projectile shows the result at work. A ball thrown at
\SI{12}{\metre\per\second} at \ang{60} above the horizontal feels only its
weight, $\vb{F} = (0, -mg)$, with $y$ upwards, so over a rise $y$ the
weight does the work $\vb{F}\cdot\Delta\vb{r} = -mgy$. At the top of the
path the velocity is horizontal (\cref{sec:mo-plane}) and equal to the
unchanging horizontal part, $12\cos\ang{60} = \SI{6}{\metre\per\second}$.
Equating the change of kinetic energy to the work and dividing by $m$,
\begin{equation*}
   \tfrac12\bigl(6^2 - 12^2\bigr) = -g\,y_{\mathrm{top}}
   \quad\Longrightarrow\quad
   y_{\mathrm{top}} = \frac{144 - 36}{2\times9.80665}
   = \SI{5.507}{\metre} ,
\end{equation*}
the height that the vertical motion alone gives, $(12\sin\ang{60})^2/2g$,
by \cref{eq:mo-v2}, since $(12\sin\ang{60})^2 = 12^2(1 - \cos^2\ang{60})
= 144 - 36$. The mass has cancelled, as it did for free fall in
\cref{sec:ne-gravity}.

Since $K$ grows as the square of the speed, doubling the speed makes it
four times larger. A car braked by a force of given size $F$, treated as a single body that
slides to rest (its spinning wheels carry a little energy of their own,
Chapter~5), stops when
that force has done work $-K$ on it. The force points against the motion,
so over a distance $d$ it does the work $-Fd$, and $-Fd = -K$ gives $d =
K/F$: at twice the speed the car needs four times the distance to stop.

\begin{selfcheck}
A body's kinetic energy is \SI{50}{\joule}. What is it when the speed is
halved? What is the kinetic energy of a \SI{2}{\kilogram} body moving at
\SI{3}{\metre\per\second}?
\end{selfcheck}
